\chapter{Design decisions and safety}

This chapter explains the technical choices and the implemented safety measures. The goal is a robust, predictable, and easily testable firmware.

\section{Modularity and responsibility}
Splitting the code into modules under \texttt{lib/} addresses three needs: isolate sensor logic, simplify reuse, and enable focused unit tests. Each module has a single responsibility and a minimal API surface.

The separation between \texttt{begin()} and \texttt{update()} distinguishes hardware configuration from runtime behavior. This reduces side effects and makes logic observable even in tests using mocks.

\section{Sensor error handling}
For analog and digital sensors, a consistent sentinel value is used: \texttt{-999.0} for read errors. This choice makes it easy to identify errors downstream without exceptions or complex state.

Each control module validates readings before actuating. If input is invalid, the module turns outputs off (fail-safe) and returns \texttt{false}. This prevents silent error propagation.

\section{Actuator protection and stability}
In the servo module, \texttt{antiSufferingServo} prevents motion beyond mechanical limits, logging a message on Serial. This avoids stalls or motor stress.

The thermostat introduces a pause, defined by \texttt{\_pauseTime}, to avoid rapid toggling when readings oscillate near the threshold. This improves reliability and actuator lifespan.

Lighting control uses a PWM range limited to 0--255. In Wokwi simulation, the value is forced to 255 when light is needed, ensuring visual consistency and reducing ambiguity in tests.

\section{Interface and communication choices}
The LCD is updated at most once per second and rotates pages every 5 seconds. This policy protects against excessive writes, avoids flicker, and preserves readability.

The error page summarizes operational issues (humidity high/low, abnormal temperature, people count out of range) with short texts. The choice is oriented toward quick diagnostics in real environments.

\section{Testing strategy}
The test suite validates edge conditions: sensor limits, actuator limits, state transitions, and safety fallback. Each test checks both logical output and side effects on mocked I/O.

Local Arduino mocks allow tests to run on a PC, speeding up feedback and reducing hardware dependence. This makes the project more maintainable and suitable for future evolution.
